\documentclass[11pt]{article}
\usepackage{amsmath, amssymb}
\usepackage[dvipsnames]{xcolor}
\usepackage[shortlabels]{enumitem}
\usepackage[margin=1in]{geometry}

\setlength{\parindent}{0pt}
\setlength{\parskip}{8pt}

%%
%% Color-box notation for the chain rule:
%%   \outside{...} draws a RED box around the entire composite expression (outer function applied to inner)
%%   \inside{...}  draws a BLUE box around the inner function
%%
\newcommand{\outside}[1]{\fcolorbox{red}{white}{#1}}
\newcommand{\inside}[1]{\fcolorbox{blue}{white}{#1}}

\title{Worksheet: The Chain Rule}
\author{Math Ma --- Single-Variable Calculus}
\date{Fall 2025}

\begin{document}
\maketitle

\section*{Identifying the Chain Rule Structure}

When differentiating a composition $f(g(x))$, we annotate the expression using
colored boxes:
\begin{itemize}
  \item A \fcolorbox{red}{white}{\textcolor{red}{red box}} indicates the \emph{outer function} applied to the inner piece.
  \item A \fcolorbox{blue}{white}{\textcolor{blue}{blue box}} indicates the \emph{inner function} $g(x)$.
\end{itemize}
Then the Chain Rule says:
\[
  \frac{d}{dx}\bigl[ f(g(x)) \bigr] = f'(g(x)) \cdot g'(x).
\]

\section*{Worked Examples}

\begin{enumerate}
\item Differentiate $h(x) = (7x^2 + 3)^5$.

\textbf{Solution.} Identify the structure:
\[
  \frac{d}{dx}\left(\outside{$\displaystyle \inside{$\displaystyle 7x^2 + 3$}^{5}$}\right)
  = 5(7x^2 + 3)^4 \cdot 14x
  = 70x(7x^2 + 3)^4.
\]

\item Differentiate $y = \sqrt{x^2 + 1}$.

\textbf{Solution.} Rewrite as $(x^2 + 1)^{1/2}$ and apply the chain rule:
\begin{align*}
  \frac{dy}{dx}
    &= \frac{d}{dx}\left(\outside{$\displaystyle
         \inside{$\displaystyle x^2 + 1$}^{1/2}$}\right) \\[6pt]
    &= \frac{1}{2}(x^2 + 1)^{-1/2} \cdot \frac{d}{dx}(x^2 + 1) \\[6pt]
    &= \frac{1}{2}(x^2 + 1)^{-1/2} \cdot 2x
     = \frac{x}{\sqrt{x^2 + 1}}.
\end{align*}

\item Differentiate $f(x) = e^{x^3}$.

\textbf{Solution.}
\begin{align*}
  f'(x)
    &= \frac{d}{dx}\left(\outside{$\displaystyle
         e^{\inside{$\displaystyle x^3$}}$}\right) \\[6pt]
    &= e^{x^3} \cdot 3x^2.
\end{align*}

\item Differentiate $y = \ln(x^2 + 5x)$.

\textbf{Solution.}
\[
  \frac{dy}{dx}
    = \frac{d}{dx}\left(\outside{$\displaystyle
       \ln\!\left(\inside{$\displaystyle x^2 + 5x$}\right)$}\right)
    = \frac{1}{x^2 + 5x} \cdot (2x + 5)
    = \frac{2x + 5}{x^2 + 5x}.
\]

\item Differentiate $g(x) = \cos^3(2x)$. \emph{This requires applying the chain rule twice!}

\textbf{Solution.} First, rewrite as $[\cos(2x)]^3$:
\begin{align*}
  g'(x) &= \frac{d}{dx}\left(
    \outside{$\displaystyle
      \inside{$\displaystyle \cos(2x)$}^{3}
    $}\right) \\[6pt]
  &= 3[\cos(2x)]^2 \cdot \frac{d}{dx}\left(
    \outside{$\displaystyle
      \cos\!\left(\inside{$\displaystyle 2x$}\right)
    $}\right) \\[6pt]
  &= 3\cos^2(2x) \cdot \bigl[-\sin(2x)\bigr] \cdot 2 \\[6pt]
  &= -6 \cos^2(2x) \sin(2x).
\end{align*}
\end{enumerate}

\section*{Practice Problems}

Use the color-box annotation to identify the outer and inner functions, then differentiate.

\begin{enumerate}[resume]
\item $f(x) = (3x^4 - x)^7$

\item $g(t) = e^{-t^2/2}$

\item $h(x) = \sin(e^x)$

\item $y = \left( \dfrac{x+1}{x-1} \right)^3$

\item $F(x) = \ln(\sin x)$

\item $G(x) = \sqrt[3]{\tan(5x)}$

\item If $r(x) = f(g(h(x)))$ with the following information:
\begin{align*}
  h(1) &= 3, & g(3) &= 2, & f'(2) &= 7, \\
  h'(1) &= 4, & g'(3) &= 5,
\end{align*}
find $r'(1)$.
\end{enumerate}

\end{document}
